\section{Related Work}
\label{sec:related}

Three lines of work bear on our study: graph-based POI embeddings that supply the substrate, multi-task learning balancers that sit on top of any joint architecture, and prior multi-task work on POI prediction tasks. We summarise each line briefly and position our contribution against it; a fourth note on cross-attention multi-task models is deferred to the methodological discussion in Section~\ref{sec:mechanism}.

\subsection{Graph-based POI embeddings}
\label{subsec:related-embeddings}

Substrate choice is the load-bearing axis our study tests. Early POI representations learn one vector per place from co-visit statistics, treating the POI as the atomic unit, and category-aware extensions inject category supervision into the same objective. Self-supervised graph contrastive methods then sharpened the representation: Deep Graph Infomax (DGI)~\cite{velickovic2019dgi} maximises mutual information between local node embeddings and a global graph summary, and Hierarchical Graph Infomax (HGI)~\cite{huang2023hgi} adapts the contrast across POI, region, and city graphs to capture multi-scale urban structure. DGI and HGI both emit \emph{POI-stable} vectors: the same place visited twice yields the same representation. The substrate we use, \emph{Check2HGI}, is a check-in-level extension of HGI's contrastive objective: contrastive views are sampled at the granularity of a single check-in, and the visit's temporal and co-visit context perturbs the view, so the same POI visited twice can yield two distinct vectors. This per-visit emission rate is the property our substrate axis tests; the recipe is described in our anonymous code release.

\subsection{Multi-task learning balancers}
\label{subsec:related-mtl}

A second body of work studies how to combine task losses or gradients once the architecture is fixed. Gradient-balancing methods rescale, project, or align per-task gradients; recent proposals along this axis include FAMO~\cite{liu2023famo}, which adapts loss weights from per-task progress, and Aligned-MTL~\cite{senushkin2023aligned}, which aligns gradient directions through an independent-component decomposition. The dense-prediction survey~\cite{vandenhende2022mtl} aggregates the line and reports that no balancer dominates across task pairs; the textbook tradeoff identified by Caruana~\cite{caruana1997multitask} and quantified by Maninis et al.~\cite{maninis2019attentive} through the $\Delta_{m}$ joint score remains the operating regime. Our setup uses a fixed static-weight aggregation with $\lambda_{\text{cat}} = 0.75$. We do not propose a new balancer, and we ablate against FAMO and Aligned-MTL in Section~\ref{sec:mechanism} to check whether a drop-in balancer recovers the next-region cost we will report.

\subsection{Prior POI multi-task work}
\label{subsec:related-poi-mtl}

A third line is concept-aligned with our task pair. The human-mobility representation model~\cite{chen2020hmrm} jointly models attribute and trajectory signals through shared encoders, and STAN~\cite{luo2021stan} addresses the next-location subproblem with spatio-temporal self-attention; we adapt STAN for our region head and add a training-fold-only additive $\alpha\!\cdot\!\log T$ prior in the spirit of GETNext~\cite{yang2022getnext}, but our head is not a faithful reproduction of any one prior model. ReHDM~\cite{li2025rehdm} extends the hypergraph line with a region-aware dual-level model and serves as our concept-aligned external next-region baseline. Two recent works on the precise category and next-POI pairing are closer still. A first study trained a hard-parameter-sharing multi-task framework over POI-stable graph embeddings (DGI), comparing against single-task baselines such as HMRM~\cite{chen2020hmrm} and MHA+PE~\cite{zeng2019mhape}, and reported only marginal joint gains, with most differences within standard deviations; the diagnosis pointed at representation mismatch, where a single shared encoder learns a compromise specialised enough for neither head~\cite{silva2025mtlnet}. A follow-up tested an embedding-side response by decomposing the monolithic graph embedding into independent spatial, temporal, and categorical sub-encoders, and recovered substantial category-side gains across three states, supporting the reading that the embedding choice was the primary factor~\cite{paiva2026courb}. Both responses kept the embedding's per-POI emission rate fixed and varied the \emph{what} axis (which features to encode); our study varies the \emph{granularity} axis instead.

Our work re-examines a default reading of multi-task learning for POI prediction: that joint training transfers signal between heads to lift both. With the substrate held fixed at the right granularity, joint cross-attention training delivers a small additional category lift at four of five Gowalla~\cite{cho2011gowalla} state splits and pays a sign-consistent cost on next-region at every state. We use a controlled five-state ablation on user-disjoint folds to decompose the joint result, separating the substrate's contribution from the architecture's.
